\section{Introducción}
\subsection{La Historia detrás del EIT}
La interacción entre la radiación electromagnética y la materia ha sido históricamente modelada, en primera aproximación, mediante sistemas semiclásicos de dos niveles. Bajo este esquema tradicional, la absorción óptica de un medio resonantemente excitado responde a perfiles simétricos continuos gobernados por ensanchamientos homogéneos (como los perfiles lorentzianos) e inhomogéneos (como el ensanchamiento Doppler y los perfiles de Voigt). Sin embargo, a medida que la física atómica ha avanzado hacia el control coherente de multinivel, surgieron fenómenos ópticos no lineales donde la coherencia de fase entre estados cuánticos altera drásticamente la respuesta macroscópica del medio.

La Transparencia Inducida Electromagnéticamente (EIT) nace de un fenómeno denominado Atrapamiento Coherente de Población (CPT), observado por primera vez por Alzetta y colaboradores en 1976 \cite{alzettaExperimentalMethodObservation1976} en un sistema atómico de tres niveles. En el CPT, dos campos ópticos acoplan simultáneamente dos subniveles del estado fundamental a un estado excitado común, induciendo una superposición coherente no radiativa conocida como estado oscuro (dark state), en la cual los electrones quedan atrapados sin posibilidad de absorber fotones.

Fue a principios de la década de 1990 cuando Harris, Field e Imamoğlu formalizaron el concepto de EIT \cite{harrisNonlinearOpticalProcesses1990}, demostrando teóricamente que era posible eliminar la absorción lineal resonante de un medio atómico mediante la aplicación de un campo de acoplamiento (coupling), preservando al mismo tiempo (e incluso amplificando) las susceptibilidades no lineales. La demostración experimental pionera fue realizada por Boller, Imamoğlu y Harris (1991) en vapores de estroncio, a la que siguieron las mediciones de dispersión y formas de línea en vapores térmicos de Rubidio desarrolladas por Xiao y Gea-Banacloche (1995) \cite{xiaoMeasurementDispersiveProperties1995}. El interés por el EIT se catapultó a finales de los 90s y principios de los 00's tras los experimentos de luz lenta realizados por el equipo de trabajo deHau et al. (1999) \cite{hauLightSpeedReduction1999} en gases ultrafríos y Kash et al. (1999) en celdas de vapor caliente, así como el posterior almacenamiento y recuperación de fotones en coherencias atómicas de espín demostrado por Phillips et al. (2001) \cite{phillipsStorageLightAtomic2001}.

Desde el punto de vista del formalismo cuántico, la EIT se manifiesta en la susceptibilidad eléctrica compleja $\chi(\omega) = \chi'(\omega) + i\chi''(\omega)$. La interferencia cuántica destructiva anula la componente imaginaria ($\text{Im}[\chi] \to 0$), abriendo una estrecha "ventana" de transparencia en el centro de la línea de absorción. Simultáneamente, la componente real ($\text{Re}[\chi]$) exhibe una dispersión anómala extremadamente abrupta en la zona de transparencia, derivando en un índice de grupo $n_g \gg 1$ responsable de la reducción drástica en la velocidad de grupo de la luz ($v_g \ll c$).
\subsection{El EIT en México: Avances y Progreso}
En México, la investigación en espectroscopía de alta resolución, coherencia atómica, átomos fríos, fotónica y materia cuántica ha sido liderada principalmente por grupos de investigación en el Instituto de Física (IF-UNAM) y el Instituto de Ciencias Nucleares (ICN-UNAM) de la Universidad Nacional Autónoma de México.

En el Instituto de Física (IF-UNAM), a través del Laboratorio de Materia Ultrafría (LMU) y el Laboratorio de Átomos Fríos y Óptica Cuántica (LAFRIOC) asociados al Laboratorio Nacional de Materia Cuántica (LANMAC), se han consolidado líneas de investigación pioneras enfocadas en el enfriamiento y atrapamiento neutro de átomos alcalinos (como Rubidio y Cesio), trampas magneto-ópticas (MOT) y el estudio de las interacciones luz-materia a bajas temperaturas. Estas plataformas experimentales han permitido explorar fenómenos de coherencia atómica tanto en regímenes de gases atómicos ultrafríos como en celdas térmicas convencionales. Asimismo, el desarrollo de infraestructura espectroscópica basada en celdas de vapor de Rubidio ha sentado las bases para el análisis detallado de la estructura hiperfina en transiciones como las líneas $D_1$ y $D_2$, donde la presencia de efectos de bombeo óptico y coherencia interna exige superar los modelos simplificados de perfiles de absorción. 

Por su parte, en el Instituto de Ciencias Nucleares (ICN-UNAM), el grupo de Óptica Cuántica y Física Atómica ha aportado avances fundamentales en el control coherente de fotones, la generación de estados entrelazados, la espectroscopía de absorción saturada y la implementación teórica y experimental de efectos no lineales acoplados a sistemas atómicos. Las colaboraciones entre ambas instituciones han fomentado la formación de recursos humanos de alto nivel e impulsado el desarrollo de instrumental óptico, caracterización térmica de vapores y sistemas de adquisición para la metrología cuántica y la metrología de frecuencias.

El estudio del EIT en celdas de vapores alcalinos no es únicamente de interés fundamental para la física cuántica de sistemas de tres niveles (configuración $\Lambda$, $V$ o escalera), sino que posee una relación directa con el modelado espectroscópico moderno. En experimentos reales con vapores atómicos calientes, la presencia de campos intensos de acoplamiento rompe la simetría de las envolventes Voigt tradicionales. Para parametrizar con precisión la forma de línea (line shape), resulta indispensable resolver la matriz de densidad ($\rho$) mediante la ecuación máster de Liouville-von Neumann e incorporar los parámetros de descoherencia ($\gamma_{12}$), frecuencias Rabi ($\Omega$) y ensanchamiento Doppler.
